\documentclass{article}
\usepackage{pathfinder-note}

\title{Forecast-Conditioned Carbon Scheduling and Nash Implementation}

\begin{document}
\maketitle

\section*{The pair}

Q proposes a quadratic payment rule to implement a convex resource allocation optimum through a strategic game. P forecasts nodal carbon intensity and uses that forecast to schedule mobile storage and data-center loads on a power network.

\section*{Connection pursued}

The proposed connection was to use P's day-ahead forecast as a carbon price in a Q-style mechanism, so that decentralized choices implement forecast-driven scheduling. The fit is limited: Q assumes convex local sets, strongly concave valuations, and componentwise aggregate capacity limits, while P's mobile-storage routing uses binary location and travel decisions; P also conserves data-center workload, whereas Q's participation benchmark uses a feasible zero allocation. \ledger{1, 3, 7, 23}

\section*{What was established}

\begin{objection}
Q's stated uniqueness certificate cannot coexist with its implementation condition. Its condition P2(i) sets \(\sum_m A_{nm}^n=0\) for each \(n\). For a vector with zero allocation coordinates and the same nonzero price vector \(q\) in every agent block, the quadratic form of Q's matrix satisfies
\[
v^\top\Psi v
=-\sum_n q^\top\Bigl(\sum_m A_{nm}^n\Bigr)q=0.
\]
Thus \(\Psi+\Psi^\top\) cannot be strictly negative definite as required by P1(ii). This refutes that joint LMI certificate, rather than implementation itself. \ledger{4}
\end{objection}

Full strategy uniqueness also fails for Q's displayed payment. With two agents, one resource, \(\alpha=1\), \(c=2\), \(\mathcal X_n=[0,1]\), and \(V_n(x)=3x-x^2/2\), every profile \((x_1,x_2,p_1,p_2)=(1,1,p,p)\) with \(0\leq p\leq2\) is an equilibrium. All have the same efficient allocation. \ledger{6}

A direct best-response argument retains an allocation-level result. Write \(S^k=\sum_n x_n^k-c^k\) and \(P^k=\sum_n p_n^k\). At an equilibrium of Q's payment rule, each price satisfies
\[
p_n^k=\left[\frac{P^k-p_n^k+S^k}{N-1}\right]_+.
\]
If one price is positive, this equation forces every price for that resource to be positive and equal, and \(S^k=0\). If all prices are zero, \(S^k\leq0\). Hence prices are common, allocations are feasible, and price times capacity slack is zero. At a common price \(p\), agent \(n\)'s allocation marginal payment is \(\alpha p\). Its best-response conditions are therefore the planner conditions with multiplier \(\lambda=\alpha p\). Given a planner primal-dual pair, the converse best-response argument supplies an equilibrium. Strong concavity makes its allocation unique, even when prices are not. \ledger{8, 12, 18}

\begin{objection}
Q's displayed rule does not always provide individual rationality. For \(N=2\), \(\alpha=c=1\), \(\mathcal X_n=[0,1]\), \(V_1(x)=3x-x^2/2\), and \(V_2(x)=2x-x^2/2\), the profile \((x_1,x_2,p_1,p_2)=(1,0,2,2)\) is an equilibrium. Q's payments are \((3,-3)\), giving agent 1 utility \(V_1(1)-3=-1/2<0\). \ledger{10, 13}
\end{objection}

The ledger gives a restricted repair. Let \(\bar p_{-n}=\sum_{m\ne n}p_m\) and \(\bar x_{-n}=\sum_{m\ne n}x_m\). Add to Q's payment the transfer
\[
h_n=
\frac{\alpha N}{(N-1)^2}
\left(\frac{\bar p_{-n}}{N-1}\right)^\top
\left(\bar x_{-n}-\frac{N-1}{N}c\right).
\]
It depends only on other agents' actions, so best responses are unchanged. At equilibrium, Q's payment is
\(\alpha[1+N/(N-1)^2]p^\top(x_n-c/N)\); the added transfer removes the excess factor, leaving \(\alpha p^\top(x_n-c/N)\). The transfers sum to zero at equilibrium. \ledger{11, 13, 18, 22}

For a frozen nonnegative forecast \(\hat e\) and optional nonnegative convex loads, also add
\[
b_n=\beta\hat e^\top x_n
-\frac{\beta}{N-1}\hat e^\top\bar x_{-n}.
\]
The second term is independent of agent \(n\)'s choice, and \(\sum_n b_n=0\) at every profile. Best responses consequently use the strongly concave valuation \(V_n(x_n)-\beta\hat e^\top x_n\). Every equilibrium implements its unique forecast-conditioned planner allocation. If \(0\in\mathcal X_n\), comparison with the feasible deviation \(x_n=0\), together with \(p,\hat e,\bar x_{-n}\geq0\), gives nonnegative equilibrium utility under the corrected payment. This is an optional-load theorem, not an implementation of P's fixed-workload or binary-routing program. \ledger{5, 7, 18, 22, 23, 25}

\section*{Carbon outcome and unresolved objections}

P calculates nodal intensity using generation, flows, and mobile-storage discharge, yet schedules against a fixed forecast. A zero-error baseline forecast need not rank the emissions caused by new load correctly: in the ledger's two-location example, a zero-carbon generator is already at capacity at location A, so an added unit uses a backup source emitting one unit; location B's available source emits one-half unit. Baseline average intensity ranks A first, while physical incremental emissions rank B first. This hypothetical example does not describe P's IEEE 33-bus simulation. \ledger{2, 15, 18, 23}

Let \(X(x)=\sum_n x_n\) and \(0\leq X(x)\leq c\). If the counterfactual intensity \(e(x)\) obeys the \emph{uniform} bound
\(\sup_{x\in\mathcal F}\|e(x)-\hat e\|_\infty\leq\varepsilon\), the repaired mechanism's forecast-optimal allocation has at most \(2\beta\varepsilon\|c\|_1\) regret for the attributed-load objective \(e(x)^\top X(x)\). When true intensity is a fixed vector \(e\), the bound sharpens to \(\beta\varepsilon\|c\|_1\). For actual generation emissions \(C(x)\), an additional uniform accounting link
\[
\sup_{x\in\mathcal F}
\left|C(x)-C_0-e(x)^\top X(x)\right|\leq\delta
\]
would give physical-welfare regret at most
\(2\beta(\varepsilon\|c\|_1+\delta)\). P's reported observational forecast errors establish neither uniform premise. \ledger{12, 18, 20, 22, 23, 26}

No stochastic-learning convergence result or mixed-integer implementation follows from this repair. The material supports a restricted allocation and payment result; it is thin on counterfactual emissions, the mandatory-workload outside option, and validation beyond that result. \ledger{7, 18, 23, 25}

\section*{Outside source and next step}

The peers checked Chen and Zhao, \href{https://arxiv.org/abs/2308.08195}{\emph{Achieving Social Optimum and Budget Balance via a Joint Electricity-Carbon Pricing Mechanism}}. Its stated social-optimum, budget-balance, and individual-rationality claims mean that a balanced carbon charge alone is not an established novelty for this pair. The ledger records no exhaustive prior-work review. \ledger{14, 19}

The smallest next step is to specify one two-location movable-load dispatch with an explicit generation-emissions function \(C(x)\), calculate its counterfactual intensity \(e(x)\), and test the two uniform bounds above. That would determine whether the restricted mechanism's attributed-load guarantee can say anything about physical emissions in that instance. \ledger{15, 23, 26}

\end{document}